\chapter{État de l'art}
\label{chap:etat-de-lart}

\section{Le modèle de Reynolds : flocking et boids}

En 1987, Craig W. Reynolds publie au SIGGRAPH un article fondateur
intitulé \emph{Flocks, herds and schools: a distributed behavioral
model}~\cite{reynolds1987flocks}. Il y propose un modèle de simulation
de groupes d'animaux (volées d'oiseaux, troupeaux, bancs de poissons)
qui repose sur trois constats~:

\begin{enumerate}
  \item Le comportement de groupe ne provient pas d'un \emph{leader}
    qui dicterait la direction commune.
  \item Chaque agent (que Reynolds appelle un \emph{boid}) prend ses
    décisions à partir d'observations \emph{purement locales}~: il
    ne voit que ses voisins dans un voisinage défini.
  \item Les règles locales sont peu nombreuses et simples, mais le
    comportement global qui en émerge est riche et crédible.
\end{enumerate}

\subsection{Les trois règles canoniques}

Reynolds définit trois règles, appliquées indépendamment puis
combinées en une accélération unique~:

\paragraph{Séparation.} Évite les collisions avec les voisins trop
proches. Mathématiquement, on accumule une force inversement
proportionnelle à la distance entre l'agent et chacun de ses voisins
proches~:
\[
  \vec{F}_{\text{sep}} \propto -\sum_{i \in V_{\text{proches}}} \frac{\vec{r}_i}{\|\vec{r}_i\|^2}
\]
où $\vec{r}_i = \vec{p}_i - \vec{p}_{\text{self}}$ est le vecteur
allant de l'agent au voisin~$i$.

\paragraph{Alignement.} Aligne sa propre vitesse sur la moyenne des
vitesses des voisins~:
\[
  \vec{F}_{\text{ali}} \propto \overline{\vec{v}}_V - \vec{v}_{\text{self}}
\]

\paragraph{Cohésion.} Tend à se rapprocher du \emph{centre de masse}
local des voisins~:
\[
  \vec{F}_{\text{coh}} \propto \overline{\vec{p}}_V - \vec{p}_{\text{self}}
\]

\subsection{Combinaison et limites}

Les trois forces sont additionnées, chacune pondérée par un
coefficient (\code{separation\_weight}, etc.), puis l'accélération
résultante est bornée par une \emph{force maximale} \code{max\_force}
afin d'éviter des changements brusques. La vitesse résultante est
elle-même bornée par \code{max\_speed}, et un seuil minimal
\code{min\_speed} maintient les agents en mouvement.

Cette construction présente plusieurs propriétés intéressantes~:

\begin{itemize}
  \item Elle est \emph{distribuée}~: chaque boid décide pour
    lui-même, sans coordination centrale.
  \item Elle est \emph{paramétrable}~: en jouant sur les rayons et
    les poids, on obtient des comportements visuellement très
    différents (banc serré, nuée diffuse, individus solitaires).
  \item Elle a une complexité naïve en $O(N^2)$ (chaque agent
    examine chaque autre), améliorable par des structures spatiales
    de type \emph{spatial hashing} ou \emph{octree}~\cite{ericson2004rtcd},
    mais cette complexité reste acceptable pour
    $N \lesssim 200$ comme dans notre cas.
\end{itemize}

\section{Comportements de pilotage : au-delà des boids}

Le modèle initial de Reynolds décrit trois forces de groupe, mais
laisse de côté les interactions inter-espèces. Dans un article
ultérieur~\cite{reynolds1999steering}, Reynolds généralise sous le
nom de \emph{steering behaviors} (comportements de pilotage) une
famille plus large de forces utiles pour les agents autonomes~: \emph{seek}
(aller vers une cible), \emph{flee} (s'éloigner), \emph{pursue}
(intercepter une cible mobile), \emph{wander} (errer aléatoirement),
\emph{obstacle avoidance}, etc. Millington et Funge approfondissent
cette taxonomie dans le contexte du jeu vidéo~\cite{millington2009ai}.

Pour notre aquarium, les forces ajoutées au-delà des trois forces
classiques sont les suivantes~:

\begin{itemize}
  \item \textbf{Fuite (\emph{flee}).} Une proie s'éloigne d'un
    prédateur visible dans son rayon de peur \code{fear\_radius},
    avec un \emph{boost} d'adrénaline \code{ADRENALINE\_BOOST} qui
    multiplie temporairement sa vitesse maximale.
  \item \textbf{Chasse (\emph{seek}).} Un prédateur se dirige vers
    sa proie la plus proche dans son rayon de chasse
    \code{hunt\_radius}, à condition que l'espèce de la proie
    figure dans son masque \code{prey\_mask}.
  \item \textbf{Évitement intra-espèce.} Une force d'écartement
    supplémentaire entre individus d'une même espèce, au-delà de
    la séparation locale, modélisée par \code{avoid\_same\_radius}.
    Cette force évite que des prédateurs solitaires (requin,
    barracuda) se collent les uns aux autres.
\end{itemize}

\section{Intégration à pas fixe vs pas variable}

Toute simulation continue échantillonnée doit choisir un schéma
d'intégration numérique. Les deux approches usuelles sont~:

\paragraph{Pas variable.} À chaque image, on calcule le délai
$\Delta t$ écoulé depuis l'image précédente, puis on intègre les
positions et vitesses sur cet intervalle. Avantage~: la simulation
suit le rythme réel, peu importe que le rendu prenne 16~ms ou
50~ms. Inconvénient majeur~: le comportement devient \emph{non
déterministe} et numériquement instable lorsque $\Delta t$ varie
fortement (typiquement à cause d'un \emph{stutter} système).

\paragraph{Pas fixe.} On choisit un pas $\Delta t$ constant (par
exemple 16~ms, soit environ 60~Hz) et on appelle la fonction de
mise à jour à intervalles réguliers. Avantage~: déterminisme
parfait, stabilité numérique, code simple. Inconvénient~: si le
rendu prend plus longtemps que $\Delta t$, le système peut
prendre du retard. Cette limitation se gère par un mécanisme de
\emph{rattrapage} (faire plusieurs ticks de simulation par image
de rendu) si nécessaire.

Pour notre projet, le compromis bénéfice/complexité est net~: le
pas fixe est implémenté trivialement avec
\code{g\_timeout\_add(TICK\_MS, \ldots)} de GLib~\cite{glib}, et la
simulation reste fluide tant que le coût de calcul d'un tick reste
inférieur à \code{TICK\_MS}.

\section{GTK4 comme moteur de rendu}

GTK4~\cite{gtk4docs} est le \emph{toolkit} graphique du projet
GNOME. Il s'agit d'une bibliothèque d'interface utilisateur de
bureau, et non d'un moteur de jeu. L'utiliser pour piloter une
simulation animée est un choix non standard qui mérite quelques
explications.

\paragraph{Arguments en faveur de GTK4.} Bibliothèque native, écrite
en C, intégration parfaite avec l'environnement Linux, gestion
automatique de la fenêtre, du plein écran, des entrées clavier,
et accès à un système de \emph{CSS} pour styler les widgets, y
compris des transformations affines (rotation, mise à l'échelle).
La présence du widget \code{GtkPicture} pour afficher une image
fixe et la possibilité de placer ces widgets en coordonnées
absolues via \code{GtkFixed} permettent d'obtenir un rendu
2D~animé sans dépendre d'OpenGL ou d'une bibliothèque tierce.

\paragraph{Arguments contre.} GTK n'est pas optimisé pour des
centaines de widgets manipulés à 60~Hz. Chaque entité de
l'aquarium étant un widget GTK distinct, on charge le système avec
$N$ recalculs de style CSS par image. En pratique, à 95~poissons,
les performances restent acceptables sur du matériel moderne, mais
ce choix limiterait la simulation à quelques centaines d'agents
au plus. Une version future utilisant Cairo ou GTK \emph{snapshot}
pour dessiner directement les sprites en une seule passe gagnerait
significativement en performance.

\section{Formats de configuration}
\label{sec:config-formats}

Le cœur du projet est de permettre la modification de paramètres
sans recompilation. Cela impose de choisir un format de
configuration. Quatre familles ont été envisagées~:

\paragraph{INI.} Très répandu, simplissime à parser, lisible.
Faiblesses~: pas de structures imbriquées standard, pas de
listes natives, types implicites (tout est chaîne).

\paragraph{JSON.} Universel, beaucoup d'outils, types stricts
(nombres, booléens, listes). Faiblesses pour un fichier
\emph{édité à la main}~: pas de commentaires, virgules
trompeuses, guillemets obligatoires partout, parenthèses
imbriquées qui rendent la structure visuelle confuse.

\paragraph{YAML.} Lisible, permet les commentaires, listes
natives. Faiblesses~: spécification très complexe (plusieurs
ambiguïtés célèbres comme \emph{Norway problem}), parseur lourd,
sensibilité à l'indentation par espaces qui pénalise les éditeurs
non-techniques.

\paragraph{TOML.} \emph{Tom's Obvious, Minimal
Language}~\cite{tomlspec}, format conçu spécifiquement pour les
fichiers de configuration. Avantages décisifs pour notre cas~:

\begin{itemize}
  \item Commentaires natifs avec~\code{\#}.
  \item Sections claires \code{[global]}, tableaux de tables
    \code{[[species]]}.
  \item Types stricts (entier, flottant, booléen, chaîne, tableau,
    date).
  \item Spécification courte et sans ambiguïté.
  \item Parseurs disponibles dans la plupart des langages, en
    particulier \code{tomlc99}~\cite{tomlc99}, une implémentation
    en C99 distribuée sur deux fichiers (\file{toml.h} et
    \file{toml.c}), facile à \emph{vendor}~iser dans un projet C
    sans dépendance externe.
\end{itemize}

Le tableau~\ref{tab:format-comparison} récapitule cette
comparaison.

\begin{table}[h!]
\centering
\footnotesize
\begin{tabular}{@{}lcccc@{}}
\toprule
\textbf{Critère} & \textbf{INI} & \textbf{JSON} & \textbf{YAML} & \textbf{TOML} \\
\midrule
Commentaires natifs    & oui & non & oui & oui \\
Listes natives         & non & oui & oui & oui \\
Types stricts          & non & oui & oui & oui \\
Tables imbriquées      & non & oui & oui & oui \\
Édition manuelle       & ++  & --  & +   & ++  \\
Parseur C minimaliste  & oui & oui & non & oui \\
Spécification courte   & oui & oui & non & oui \\
\bottomrule
\end{tabular}
\caption{Comparaison des formats de configuration courants.}
\label{tab:format-comparison}
\end{table}

\section{Synthèse}

Le projet se positionne donc à l'intersection de trois traditions~:
le modèle de Reynolds pour le comportement, l'écosystème de
développement GTK pour le rendu, et la philosophie \emph{Unix} de
configuration par fichier texte (avec chemins de recherche en
\file{./}, \file{\textasciitilde/.config/} et \file{/etc/}) pour
l'extensibilité. Chacun de ces piliers est suffisamment stable et
documenté pour servir de fondation~; la valeur ajoutée du projet
réside dans leur intégration cohérente.
